\clearpage
\section{Maxwell model in dividing tissues}
\begin{colbox}
    Consider a dividing tissue, where the cellular density satisfies
    \begin{equation}
        \del_t\rho = -\del_\alpha\ins{\rho v_\alpha} - K(\rho)
    \end{equation}
    the reaction term on RHS describes division and death. Suppose that close to the homeostatic (steady-state) density \(\rho_h\), the reaction is given by \(K(\rho)\approx\frac{1}{\tau}\ins{\rho-\rho_h}\) with \(\tau\) a typical density relaxation time due to division and dreath. In addition, assume that the pressure around the homeostatic pressure \(P_h\) satisfies to linear order \(P-P_h=K(\rho-\rho_h)\). Denote \(\delta\sigma = P_h-P\). Linearize the cellular density equation around the homeostatic state and show that \(\delta\sigma\) behaves like a Maxwell model. What is the viscosity in this case? Note that this is a very simplified model that neglects, for example, cell-cell interactions.
\end{colbox}
Since we are near the homeostatic state we can insert the relations given
\begin{equation}
    \del_t\rho = -\del_\alpha\ins{\rho v_\alpha} - \frac{1}{\tau}\ins{\rho-\rho_h}
\end{equation}
We are looking at states near the homeostatic, so we can replace \(\rho = \rho_h +\delta\rho(t)\) therefore
\begin{align}
    \del_t\ins{\delta\rho(t)} &= -\del_\alpha \ins{\ins{\rho_h+\delta\rho}v_\alpha}-\frac{\delta\rho(t)}{\tau} \\
    &= -\rho_h\del_\alpha v_\alpha - \del_\alpha\ins{\delta\rho(t) v_\alpha}-\frac{\delta\rho(t)}{\tau} \label{28}
\end{align}
The 2nd elementin equation~\ref{28} contains two very small elements since we're near the homeostatic state so we can consider it negligible
\begin{equation}
    \del_t\delta\rho(t) = -\rho_h\del_\alpha v_\alpha - \frac{\delta\rho(t)}{\tau}
\end{equation}
next we notice that
\begin{equation}
    \delta\rho(t) = \frac{P-P_h}{K} = -\frac{\delta\sigma}{K}
\end{equation}
inserting this
\begin{equation}
    \frac{1}{K}\del_t\delta\sigma = \rho_h \del_\alpha v_\alpha - \frac{\delta\sigma}{\tau K}
\end{equation}
rearranging
\begin{equation}
    (1+\tau\del_t)\delta\sigma = \tau\rho_h K\del_\alpha v_\alpha
\end{equation}
Which is the same structure as the Maxwell model. We can identify \(\tau\rho_h K\) as the viscosity.